\documentclass[UTF8]{ctexart}
\usepackage{amsmath, amssymb, geometry}
\usepackage{xcolor}
\usepackage{enumitem}
\usepackage{tcolorbox}
\usepackage{cases}
\usepackage{fancyhdr}
\geometry{a4paper, margin=1in}

% 定义红色环境
\newenvironment{redenv}{\color{red}}{}

% 定义详细解析环境
\newenvironment{detailedanalysis}{%
    \color{red}
    \begin{enumerate}[label=\textbf{第\arabic*步：}, leftmargin=*, align=left, itemsep=1.5ex]
}{%
    \end{enumerate}
}

% 定义概念框
\newenvironment{conceptbox}{%
    \begin{tcolorbox}[colback=red!5, colframe=red!50!black, title=概念解析, fonttitle=\bfseries]
}{%
    \end{tcolorbox}
}

% 定义公式推导环境
\newenvironment{formuladerivation}{%
    \begin{enumerate}[label={\textbf{公式\arabic*：}}, leftmargin=*, align=left]
}{%
    \end{enumerate}
}

% 定义题目和解析的分隔线
\newcommand{\problemrule}{\noindent\rule{\textwidth}{1.5pt}\vspace{-0.8em}}
\newcommand{\analysisrule}{\noindent\rule{\textwidth}{0.8pt}\vspace{-0.8em}}

% 宏定义：方便排版选择题选项
\newcommand{\choicesFour}[4]{
    \begin{enumerate}[label=\Alph*., itemsep=0pt, parsep=0pt, topsep=5pt, labelsep=0.5em]
        \begin{minipage}{0.22\linewidth} \item #1 \end{minipage}
        \begin{minipage}{0.22\linewidth} \item #2 \end{minipage}
        \begin{minipage}{0.22\linewidth} \item #3 \end{minipage}
        \begin{minipage}{0.22\linewidth} \item #4 \end{minipage}
    \end{enumerate}
}
\newcommand{\choicesTwo}[4]{
    \begin{enumerate}[label=\Alph*., itemsep=0pt, parsep=0pt, topsep=5pt, labelsep=0.5em]
        \begin{minipage}{0.45\linewidth} \item #1 \end{minipage}
        \begin{minipage}{0.45\linewidth} \item #2 \end{minipage}
        \begin{minipage}{0.45\linewidth} \item #3 \end{minipage}
        \begin{minipage}{0.45\linewidth} \item #4 \end{minipage}
    \end{enumerate}
}
\newcommand{\choices}[4]{
    \begin{enumerate}[label=\Alph*., itemsep=0pt, parsep=0pt, topsep=5pt, labelsep=0.5em]
        \item #1
        \item #2
        \item #3
        \item #4
    \end{enumerate}
}

\newcommand{\blank}[1]{\underline{\hspace{#1}}}

\begin{document}

\section*{第四章 导数与微分（提高）}

% ========== 第1题 ==========
\problemrule
\section*{【题目1】导数的定义}
设 $ f(0)=0 $ 且 $ \lim_{x\to0}\frac{f(x)}{x} $ 存在，则 $ \lim_{x\to0}\frac{f(x)}{x}= $ \blank{1cm}。
\choicesFour{$ f'(x) $}{$ f'(0) $}{$ f(0) $}{$ \frac{1}{2}f'\left(\frac{1}{2}\right) $}

\begin{redenv}
\analysisrule
\subsection*{逐步解析}

\begin{detailedanalysis}
\item \textbf{问题本质分析}
本题考查导数在 $x=0$ 处的定义。

\item \textbf{详细求解过程}
导数定义：$f'(x_0) = \lim_{x \to x_0} \frac{f(x) - f(x_0)}{x - x_0}$。
取 $x_0 = 0$，并已知 $f(0)=0$。
$$ f'(0) = \lim_{x \to 0} \frac{f(x) - f(0)}{x - 0} = \lim_{x \to 0} \frac{f(x)}{x} $$
题目已知极限存在，故该极限值即为 $f'(0)$。

\item \textbf{最终答案}
\textbf{答案}：B
\end{detailedanalysis}
\end{redenv}

% ========== 第2题 ==========
\problemrule
\section*{【题目2】左右导数与可导}
函数 $ y=f(x) $ 在点 $ x_{0} $ 处的左导数 $ f'_{-}(x_{0}) $ 和右导数 $ f'_{+}(x_{0}) $ 都存在,是 $ f(x) $ 在 $ x_{0} $ 可导的 \blank{1cm}。
\choicesTwo
{充分必要条件}
{充分但非必要条件}
{必要但非充分条件}
{既非充分又非必要条件}

\begin{redenv}
\analysisrule
\subsection*{逐步解析}

\begin{detailedanalysis}
\item \textbf{详细求解过程}
可导的充要条件是：左导数存在、右导数存在，且两者\textbf{相等}。
题目只说了“都存在”，未说“相等”。
例如 $f(x) = |x|$ 在 $x=0$ 处，左导数-1，右导数1，都存在但不相等，故不可导。
所以，“都存在”是可导的必要条件（因为可导必存在），但不是充分条件（因为存在不一定相等）。

\item \textbf{最终答案}
\textbf{答案}：C
\end{detailedanalysis}
\end{redenv}

% ========== 第3题 ==========
\problemrule
\section*{【题目3】绝对值函数的导数}
函数 $ f(x)=|\sin x| $ 在 $ x=0 $ 处 \blank{1cm}。
\choicesTwo
{可导}
{连续但不可导}
{不连续}
{极限不存在}

\begin{redenv}
\analysisrule
\subsection*{逐步解析}

\begin{detailedanalysis}
\item \textbf{详细求解过程}
\begin{enumerate}[label=(\arabic*)]
\item \textbf{连续性}：$|\sin x|$ 是初等函数复合，显然连续。
\item \textbf{可导性}：考虑左右导数。
$$ f'_+(0) = \lim_{x \to 0^+} \frac{|\sin x| - 0}{x} = \lim_{x \to 0^+} \frac{\sin x}{x} = 1 $$
$$ f'_-(0) = \lim_{x \to 0^-} \frac{|\sin x| - 0}{x} = \lim_{x \to 0^-} \frac{-\sin x}{x} = -1 $$
\item \textbf{结论}：左导数 $\neq$ 右导数，故不可导。
\end{enumerate}

\item \textbf{最终答案}
\textbf{答案}：B
\end{detailedanalysis}
\end{redenv}

% ========== 第4题 ==========
\problemrule
\section*{【题目4】微分计算}
设 $ y = \tan^{2}x $ , 则 $ dy = $ \blank{1cm}。
\choicesTwo
{$ 2\tan x \sec^{2} x dx $}
{$ 2\sin x \cos^{2} x dx $}
{$ 2\sec x \tan^{2} dx $}
{$ 2\cos x \sin^{2} x dx $}

\begin{redenv}
\analysisrule
\subsection*{逐步解析}

\begin{detailedanalysis}
\item \textbf{详细求解过程}
利用复合函数求导法则（链式法则）：
$$ y' = (\tan^2 x)' = 2\tan x \cdot (\tan x)' = 2\tan x \cdot \sec^2 x $$
所以微分 $dy = y' dx = 2\tan x \sec^2 x dx$。

\item \textbf{最终答案}
\textbf{答案}：A
\end{detailedanalysis}
\end{redenv}

% ========== 第5题 ==========
\problemrule
\section*{【题目5】函数奇偶性与导数}
设 $ F(x)=f(x)+f(-x) $ ，且 $ f'(x) $ 存在，则 $ F'(x) $ 是 \blank{1cm}。
\choicesTwo
{奇函数}
{偶函数}
{非奇非偶的函数}
{无法判定其奇偶性}

\begin{redenv}
\analysisrule
\subsection*{逐步解析}

\begin{detailedanalysis}
\item \textbf{详细求解过程}
\begin{enumerate}[label=(\arabic*)]
\item \textbf{分析 F(x)}：$F(-x) = f(-x) + f(x) = F(x)$，所以 $F(x)$ 是偶函数。
\item \textbf{性质应用}：偶函数的导数是奇函数。
\item \textbf{直接求导验证}：
$$ F'(x) = f'(x) + f'(-x) \cdot (-1) = f'(x) - f'(-x) $$
计算 $F'(-x)$：
$$ F'(-x) = f'(-x) - f'(-(-x)) = f'(-x) - f'(x) = -(f'(x) - f'(-x)) = -F'(x) $$
所以 $F'(x)$ 是奇函数。
\end{enumerate}

\item \textbf{最终答案}
\textbf{答案}：A
\end{detailedanalysis}
\end{redenv}

% ========== 第6题 ==========
\problemrule
\section*{【题目6】周期函数的导数}
设周期函数 $ f(x) $ 在 $ (-\infty, \infty) $ 周期为 3, 则曲线 $ y = f(x) $ 在点 $ (4, f(4)) $ 的切线斜率为 \blank{2cm}。

\begin{redenv}
\analysisrule
\subsection*{逐步解析}

\begin{detailedanalysis}
\item \textbf{问题本质分析}
此题题目信息似乎不全（未给出具体的函数表达式或导数值）。结合常见题型推测，可能隐含了其他条件（例如是奇函数且在某点及周期点的情况），或者题目意在考查周期性关系。
但如果题目本意是求符号表达式或 specific value，必须已知 $f'(4)$ 或 $f'(1)$。
若题目意在考察 $f'(4) = f'(1)$，则答案为 $f'(1)$。
\textbf{自我修正}：如果是填空题且没有具体值，通常答案是一个具体数字。可能题目漏掉了 "$f(x)$是偶函数" 或者 "$f'(1)=...$" 的条件。这类题目如果是 "在 ... 切线斜率为 \_\_"，通常 imply 那一点是极值点或者是0。
\textbf{假设}：这是一个基于对称性的题目？
暂无法给出确切数值答案。根据常见套路，若 $f(x)$ 是可导周期函数，周期为3。如果问切线斜率，通常要用到 $f'(x)$ 的周期性 $f'(x+3) = f'(x)$，所以 $k = f'(4) = f'(1)$。
\textbf{注}：作为解析生成，我会标注出关系。

\item \textbf{详细求解过程}
因为 $f(x)$ 周期为 3，即 $f(x+3) = f(x)$。
两边求导：$f'(x+3) = f'(x)$。即导函数也以 3 为周期。
曲线在 $(4, f(4))$ 处的切线斜率 $k = f'(4)$。
根据周期性，$f'(4) = f'(1 + 3) = f'(1)$。
由于题目信息有限，答案应为 $f'(1)$（或者题目有缺漏，需检查原题）。
\end{detailedanalysis}
\end{redenv}

% ========== 第7题 ==========
\problemrule
\section*{【题目7】导数计算}
若 $ y = e^{x}(\sin x + \cos x) $ ，则 $ \frac{dy}{dx} = $ \blank{2cm}。

\begin{redenv}
\analysisrule
\subsection*{逐步解析}

\begin{detailedanalysis}
\item \textbf{详细求解过程}
利用乘法法则：
$$ y' = (e^x)'(\sin x + \cos x) + e^x (\sin x + \cos x)' $$
$$ = e^x(\sin x + \cos x) + e^x (\cos x - \sin x) $$
$$ = e^x (\sin x + \cos x + \cos x - \sin x) $$
$$ = e^x (2\cos x) = 2e^x \cos x $$

\item \textbf{最终答案}
\textbf{答案}：$ 2e^x \cos x $
\end{detailedanalysis}
\end{redenv}

% ========== 第8题 ==========
\problemrule
\section*{【题目8】隐函数求导与法线}
曲线 $ xy + \ln y - 1 = 0 $ 在点 $(1,1)$ 处的法线方程是 \blank{2cm}。

\begin{redenv}
\analysisrule
\subsection*{逐步解析}

\begin{detailedanalysis}
\item \textbf{详细求解过程}
\begin{enumerate}[label=(\arabic*)]
\item \textbf{隐函数求导}：方程两边对 $x$ 求导。
$$ 1 \cdot y + x \cdot y' + \frac{1}{y} \cdot y' = 0 $$
$$ y' (x + \frac{1}{y}) = -y $$
$$ y' = \frac{-y}{x + \frac{1}{y}} = \frac{-y^2}{xy + 1} $$
\item \textbf{求切线斜率}：代入点 $(1,1)$。
$$ k_{T} = y'|_{(1,1)} = \frac{-1^2}{1\cdot 1 + 1} = -\frac{1}{2} $$
\item \textbf{求法线斜率}：$k_n = -\frac{1}{k_T} = 2$。
\item \textbf{写方程}：点斜式 $y - 1 = 2(x - 1)$。
整理得：$2x - y - 1 = 0$。
\end{enumerate}

\item \textbf{最终答案}
\textbf{答案}：$ 2x - y - 1 = 0 $
\end{detailedanalysis}
\end{redenv}

% ========== 第9题 ==========
\problemrule
\section*{【题目9】参数方程求导}
已知 $ \begin{cases} x = a(\sin t - t\cos t) \\ y = a(\cos t + t\sin t) \end{cases} $ ，则 $ \left.\frac{dy}{dx}\right|_{t=\frac{3}{4}\pi} = $ \blank{2cm}。

\begin{redenv}
\analysisrule
\subsection*{逐步解析}

\begin{detailedanalysis}
\item \textbf{详细求解过程}
\begin{enumerate}[label=(\arabic*)]
\item \textbf{计算 $dx/dt$ 和 $dy/dt$}：
$$ \frac{dx}{dt} = a(\cos t - (\cos t - t\sin t)) = at\sin t $$
$$ \frac{dy}{dt} = a(-\sin t + (\sin t + t\cos t)) = at\cos t $$
\item \textbf{计算 $dy/dx$}：
$$ \frac{dy}{dx} = \frac{dy/dt}{dx/dt} = \frac{at\cos t}{at\sin t} = \cot t $$
\item \textbf{代入数值}：
当 $t = \frac{3}{4}\pi$ 时，
$$ \frac{dy}{dx} = \cot \frac{3}{4}\pi = -1 $$
\end{enumerate}

\item \textbf{最终答案}
\textbf{答案}：-1
\end{detailedanalysis}
\end{redenv}

% ========== 第10题 ==========
\problemrule
\section*{【题目10】高阶导数}
若 $ f(x) = x^{2} \ln x $ , 则 $ f''(2) = $ \blank{2cm}。

\begin{redenv}
\analysisrule
\subsection*{逐步解析}

\begin{detailedanalysis}
\item \textbf{详细求解过程}
\begin{enumerate}[label=(\arabic*)]
\item \textbf{一阶导数}：
$$ f'(x) = 2x \ln x + x^2 \cdot \frac{1}{x} = 2x \ln x + x $$
\item \textbf{二阶导数}：
$$ f''(x) = (2 \ln x + 2x \cdot \frac{1}{x}) + 1 = 2 \ln x + 2 + 1 = 2 \ln x + 3 $$
\item \textbf{求值}：
$$ f''(2) = 2 \ln 2 + 3 $$
\end{enumerate}

\item \textbf{最终答案}
\textbf{答案}：$ 2\ln 2 + 3 $
\end{detailedanalysis}
\end{redenv}

% ========== 第11题 ==========
\problemrule
\section*{【题目11】二阶隐函数导数}
对方程 $ e^{y} + xy = e^{2} $ 所确定的函数 $ y = y(x) $ 求 $ \frac{d^{2}y}{dx^{2}} $。

\begin{redenv}
\analysisrule
\subsection*{逐步解析}

\begin{detailedanalysis}
\item \textbf{详细求解过程}
\begin{enumerate}[label=(\arabic*)]
\item \textbf{求一阶导}：
方程两边对 $x$ 求导：
$$ e^y y' + y + xy' = 0 $$
$$ y'(e^y + x) = -y \implies y' = \frac{-y}{e^y + x} $$
\item \textbf{求二阶导}：
对 $-y = y'(e^y + x)$ 两边再求导（或者对分式求导）：
$$ y'' = \left( \frac{-y}{e^y + x} \right)' = - \frac{y'(e^y + x) - y(e^y y' + 1)}{(e^y + x)^2} $$
将 $y'(e^y+x) = -y$ 代入分子简化：
$$ \text{分子} = -y - y(e^y y' + 1) = -y - ye^y y' - y = -2y - y e^y y' $$
$$ y'' = - \frac{-2y - y e^y y'}{(e^y + x)^2} = \frac{y(2 + e^y y')}{(e^y + x)^2} $$
将 $y'$ 表达式代入：
$$ y'' = \frac{y [2 + e^y (\frac{-y}{e^y + x})]}{(e^y + x)^2} = \frac{y [2(e^y+x) - y e^y]}{(e^y+x)^3} $$
$$ = \frac{2y(e^y+x) - y^2 e^y}{(e^y+x)^3} $$
\end{enumerate}

\item \textbf{最终答案}
\textbf{答案}：$ \frac{2y(e^y+x) - y^2 e^y}{(e^y+x)^3} $
\end{detailedanalysis}
\end{redenv}

% ========== 第12题 ==========
\problemrule
\section*{【题目12】分段函数可导性}
设函数 $ f(x)=\begin{cases}ax+b, & x<1\\ x^{2}, & x\geq1\end{cases} $ 在 $x=1$ 处可导，求 $a,b$ 的值。

\begin{redenv}
\analysisrule
\subsection*{逐步解析}

\begin{detailedanalysis}
\item \textbf{详细求解过程}
可导必连续。
\begin{enumerate}[label=(\arabic*)]
\item \textbf{连续性条件}：
$\lim_{x \to 1^-} (ax+b) = a+b$。
$f(1) = 1^2 = 1$。
所以 $a+b=1$。
\item \textbf{可导性条件}（左右导数相等）：
左导数：$(ax+b)' = a$。
右导数：$(x^2)'|_{x=1} = 2x|_{x=1} = 2$。
所以 $a=2$。
\item \textbf{计算b}：
代入 $a+b=1$，得 $2+b=1 \implies b=-1$。
\end{enumerate}

\item \textbf{最终答案}
\textbf{答案}：$ a=2, b=-1 $
\end{detailedanalysis}
\end{redenv}

% ========== 第13题 ==========
\problemrule
\section*{【题目13】复合求导}
求函数 $ y = \ln(x + \sqrt{x^{2} + a^{2}}) $ 的导数。

\begin{redenv}
\analysisrule
\subsection*{逐步解析}

\begin{detailedanalysis}
\item \textbf{详细求解过程}
$$ y' = \frac{1}{x + \sqrt{x^2+a^2}} \cdot (x + \sqrt{x^2+a^2})' $$
$$ = \frac{1}{x + \sqrt{x^2+a^2}} \cdot (1 + \frac{1}{2\sqrt{x^2+a^2}} \cdot 2x) $$
$$ = \frac{1}{x + \sqrt{x^2+a^2}} \cdot (1 + \frac{x}{\sqrt{x^2+a^2}}) $$
通分括号内：
$$ = \frac{1}{x + \sqrt{x^2+a^2}} \cdot \frac{\sqrt{x^2+a^2} + x}{\sqrt{x^2+a^2}} $$
约分：
$$ = \frac{1}{\sqrt{x^2+a^2}} $$

\item \textbf{最终答案}
\textbf{答案}：$ \frac{1}{\sqrt{x^2+a^2}} $
\end{detailedanalysis}
\end{redenv}

% ========== 第14题 ==========
\problemrule
\section*{【题目14】复杂隐函数二阶导}
对方程 $ \ln(\sqrt{x^{2}+y^{2}})=\arctan\frac{y}{x} $ 所确定的函数 $ y=y(x) $ 求 $ \frac{d^{2}y}{dx^{2}} $。

\begin{redenv}
\analysisrule
\subsection*{逐步解析}

\begin{detailedanalysis}
\item \textbf{详细求解过程}
\begin{enumerate}[label=(\arabic*)]
\item \textbf{化简方程}：
左边 $= \frac{1}{2} \ln(x^2+y^2)$。
两边对 $x$ 求导：
$$ \frac{1}{2} \frac{2x + 2yy'}{x^2+y^2} = \frac{1}{1+(\frac{y}{x})^2} \cdot \frac{xy' - y}{x^2} $$
$$ \frac{x+yy'}{x^2+y^2} = \frac{x^2}{x^2+y^2} \cdot \frac{xy'-y}{x^2} = \frac{xy'-y}{x^2+y^2} $$
$$ x + yy' = xy' - y $$
$$ x + y = xy' - yy' = y'(x-y) $$
$$ y' = \frac{x+y}{x-y} $$
\item \textbf{求二阶导}：
$$ y'' = \left( \frac{x+y}{x-y} \right)' = \frac{(1+y')(x-y) - (x+y)(1-y')}{(x-y)^2} $$
展开分子：
$$ (x-y + xy' - yy') - (x - xy' + y - yy') = -2y + 2xy' = 2(xy' - y) $$
将 $y'(x-y) = x+y$ 中的 $xy' = x+y+yy'$ 代入？不如直接代入 $y'$ 值。
或者利用 $xy' - y = x$ (从 $x+yy'=xy'-y$ 得出 $x+y=y'(x-y)$... 等等，上面推导是 $x+yy' = xy'-y \implies x+y = y'(x-y)$。正确。)
回看 $x+yy' = xy'-y \implies x+y = y'(x-y)$。
我们有 $xy' - y = x + yy'$。所以分子 $2(xy'-y) = 2(x+yy')$。这个好像没简化。
直接算：
分子 $= 2(xy' - y)$。代入 $y' = \frac{x+y}{x-y}$。
$$ 2(x \frac{x+y}{x-y} - y) = 2 \frac{x^2+xy - xy + y^2}{x-y} = \frac{2(x^2+y^2)}{x-y} $$
所以
$$ y'' = \frac{2(x^2+y^2)}{(x-y)^3} $$
\end{enumerate}

\item \textbf{最终答案}
\textbf{答案}：$ \frac{2(x^2+y^2)}{(x-y)^3} $
\end{detailedanalysis}
\end{redenv}

% ========== 第15题 ==========
\problemrule
\section*{【题目15】导数化简}
求函数 $ y = \frac{x}{2} \sqrt{a^{2} - x^{2}} + \frac{a^{2}}{2} \arcsin \frac{x}{a} \quad (a > 0) $ 的导数。

\begin{redenv}
\analysisrule
\subsection*{逐步解析}

\begin{detailedanalysis}
\item \textbf{详细求解过程}
\begin{enumerate}[label=(\arabic*)]
\item 第一项求导：
$$ (\frac{x}{2} \sqrt{a^2-x^2})' = \frac{1}{2} \sqrt{a^2-x^2} + \frac{x}{2} \frac{-2x}{2\sqrt{a^2-x^2}} = \frac{\sqrt{a^2-x^2}}{2} - \frac{x^2}{2\sqrt{a^2-x^2}} $$
通分：
$$ = \frac{a^2-x^2 - x^2}{2\sqrt{a^2-x^2}} = \frac{a^2-2x^2}{2\sqrt{a^2-x^2}} $$
\item 第二项求导：
$$ (\frac{a^2}{2} \arcsin \frac{x}{a})' = \frac{a^2}{2} \frac{1}{\sqrt{1-(\frac{x}{a})^2}} \cdot \frac{1}{a} = \frac{a}{2} \frac{1}{\frac{\sqrt{a^2-x^2}}{a}} = \frac{a^2}{2\sqrt{a^2-x^2}} $$
\item 相加：
$$ y' = \frac{a^2-2x^2}{2\sqrt{a^2-x^2}} + \frac{a^2}{2\sqrt{a^2-x^2}} = \frac{2a^2-2x^2}{2\sqrt{a^2-x^2}} = \frac{2(a^2-x^2)}{2\sqrt{a^2-x^2}} = \sqrt{a^2-x^2} $$
\end{enumerate}

\item \textbf{最终答案}
\textbf{答案}：$ \sqrt{a^2-x^2} $
\end{detailedanalysis}
\end{redenv}

\end{document}
